\documentclass[blue]{beamer}

\mode<presentation> {
  \setbeamertemplate{background canvas}[vertical shading][bottom=red!10,top=blue!10]

  \usetheme{Boadilla}
  \usefonttheme[onlysmall]{serif}
}


\usepackage{pgf,pgfarrows,pgfnodes,pgfautomata,pgfheaps,pgfshade}
\usepackage{amsmath,amssymb,bm}
\usepackage{colortbl}
\usepackage[english]{babel}


\setbeamercovered{dynamic}


\def\R{\mathbb{R}}
\def\Q{\mathbb{Q}}
\def\N{\mathbb{N}}
\def\E{\mathbb{E}}
\newcommand{\ps}[2]{ \langle #1 , #2 \rangle }
\def\1{\mathbf{1}}
\newcommand{\nor}[1]{\left\| #1 \right\|}
\def\leq{\leqslant}
\def\geq{\geqslant}
\def\et{\quad \textrm{and} \quad}
\def\ou{\quad \textrm{or} \quad}
\def\boite{\hskip 0.5mm {\Box} \hskip 0.5mm}


\DeclareMathOperator\PPE{PPE}
\DeclareMathOperator\vNM{vNM}
\DeclareMathOperator\argmax{argmax}
\DeclareMathOperator\Prob{Prob}
\DeclareMathOperator\supp{supp}
\DeclareMathOperator\Eq{Eq}
\DeclareMathOperator\inte{int}
\DeclareMathOperator\co{co}
\DeclareMathOperator\cl{cl}
\DeclareMathOperator\Dirac{Dirac}
\DeclareMathOperator\F{F}
\DeclareMathOperator\As{As}
\DeclareMathOperator\Ir{Ir}
\DeclareMathOperator\MRS{MRS}
\DeclareMathOperator\Range{Im}
\DeclareMathOperator\DIV{DIV}

%------------------------------ theoren styles

\theoremstyle{definition}
\newtheorem*{assumption}{Assumption}
\newtheorem*{conjecture}{Conjecture}

\theoremstyle{example}
\newtheorem*{proposition}{Proposition}


%-----------------------------------------

\title[Corporate Finance]{Theory of Corporate Finance}
\author[Victor Filipe]{V. Filipe Martins-da-Rocha}
\institute[FGV EESP]{Escola de Economia de S\~ao Paulo}
\date[April-June, 2026]{Part 1.b. Stock Market Equilibrium: Capital Structure of the Firm}
\subject{Economic Theory}

\begin{document}

\frame{\titlepage}

%------------------ outline

\begin{frame}\frametitle{Outline: the financial structure of the firm}
    \begin{itemize}
    \item Irrelevance results: the Modigliani--Miller's theorems
    \item Taxation
    \item Limited liability and bankruptcy
    \item Irrelevance of the payout policy
    \end{itemize}
\end{frame}



%------------------------------------------------

\begin{frame}\frametitle{The financial structure of the firm: literature}

\begin{thebibliography}{8}


  \beamertemplatearticlebibitems

    \bibitem{ModiglinaiMiller58}
    Modigliani, F. and Miller, M.
    \newblock The cost of capital, corporation finance, and the theory of investment.
    \newblock {\em American Economic Review} 48, 1958

    \bibitem{Stiglitz69}
    Stiglitz, J.
    \newblock A re-examination of the Modigliani--Miller theorem.
    \newblock {\em American Economic Review} 59, 1969

    \bibitem{Stiglitz74}
    Stiglitz, J.
    \newblock On the irrelevance of corporation financial policy.
    \newblock {\em American Economic Review} 64, 1974

%    \bibitem{DeMarzo88}
%   DeMarzo, P. M.
%    \newblock An extension of the Modigliani--Miller theorem to stochastic economies with incomplete markets and interdependent securities.
%    \newblock {\em Journal of Economic Theory} 45, 1988


    \end{thebibliography}

\end{frame}

%------------------------------------------------

\begin{frame}\frametitle{The financial structure of the firm: literature}

\begin{thebibliography}{8}


  \beamertemplatearticlebibitems
  
  \bibitem{DeMarzo88}
   DeMarzo, P. M.
    \newblock An extension of the Modigliani--Miller theorem to stochastic economies with incomplete markets and interdependent securities.
    \newblock {\em Journal of Economic Theory} 45, 1988
    
    \bibitem{DeMarzoDuffie91}
   DeMarzo, P. M. and Duffie, D.
    \newblock Corporate financial hedging with proprietary information.
    \newblock {\em Journal of Economic Theory} 53, 1991

    \bibitem{Purnanandam08}
    Purnanandam, A.
    \newblock Financial distress and corporate risk management: Theory and evidence
    \newblock {\em Journal of Financial Economics} 87, 2008


    \end{thebibliography}

\end{frame}


%------------------------------------------------

\begin{frame}\frametitle{The financial structure of the firm}

\begin{itemize}
\item From now on, we fix a production plan $y$ and we focus on the relevance of the financial plan $\beta$
\item Assume the objective of the manager is to maximize the value
\[
V = E + D + \textcolor{blue}{y_0}
\]
\item Recall that
    \begin{itemize}
    \item the (present value of the) debt is given by $D = q \cdot \beta$
    \item the equity price $E$ (and the security price $q$) is determined at equilibrium and may depend on $\beta$
    \end{itemize}
\item The manager maximizes the function $\beta \mapsto E(\beta) + q \cdot \beta$
\item Assume for simplicity that $y_0=0$
\item There is a natural question: \textcolor{red}{Is there an optimal choice $\beta^\star$?}
\item The answer is: \textcolor{red}{the financial policy of the firm is irrelevant}
\end{itemize}

\end{frame}

%------------------------------------------------

\begin{frame}\frametitle{A simple case}

\begin{itemize}
\item Consider that there is a riskless asset $a_0$ promising $1$, i.e.,
\[
\forall s\in S, \quad \kappa_s(a_0) =1
\]
We denote by $r_0$ the risk free interest rate, i.e.,
\[
q(a_0) = \frac{1}{1+r_0}
\]
\item Assume that the firm can only issue (sell) the riskless bond $a_0$, i.e.,
\[
\exists Z \geq 0, \quad \beta = Z \1_{a_0}
\]
In other words, the firm commits to deliver at $t=1$ the amount $Z$ independent of the state of nature
\end{itemize}

\end{frame}

%------------------------------------------------

\begin{frame}\frametitle{A simple case}

 The value of the firm is equity plus debt, i.e.,
\[
V = E + D \quad \textrm{where} \quad D = \frac{1}{1+r_0} Z
\]
\begin{itemize}
\item the firm is said levered if $D \neq 0$
\item the firm is said unlevered if $D =0$
\item the debt-equity ratio is defined by
\[
\frac{D}{E}
\]
\item a specific choice of $\beta$ leads to a specific debt equity ration $D(\beta)/E(\beta)$
\end{itemize}

\end{frame}

%------------------------------------------------

\begin{frame}\frametitle{The conventional approach before 1958}

Assume the firm starts without debt
    \begin{itemize}
    \item the dividend of the firm $\delta_1= (y_s)_{s\in S}$  is risky
    \item debt is not risky, issuing debt may attract investors and increase the firm's value
        \begin{itemize}
        \item no debt: investors can buy the risky cash flow $(y_s)_{s\in S}$
        \item with debt: investors can still buy a risky cash flow $(y_s - Z)_{s\in S}$ but now they can also buy a riskless cash flow $Z \1_S$
        \end{itemize}
    \item a low level of debt reduces the contingencies where the firm will be unable to meet its obligations
    \item however, for a high level of debt, the possibility for bankruptcy becomes non-negligible and the benefits of leverage (issuing non-risky asset)
disappears
    \item it seems that there is an optimal level of debt
    \end{itemize}

\end{frame}


%------------------------------------------------

\begin{frame}\frametitle{Weak irrelevance result: Modigliani--Miller (1958)}

\begin{theorem}
Consider a stock market equilibrium
\[
\left\{ (y(j),\beta(j))_{j\in J}, (q,E), (c^i,\theta^i,\alpha^i)_{i\in I}\right\}
\]
Assume that
    \begin{itemize}
    \item the two firms $j_1$ and $j_2$ have the same production plan, i.e., $y(j_1)=y(j_2)$
    \item but may have different financial policies $\beta(j_1) \neq \beta(j_2)$
    \end{itemize}

Then the two firms have the same value, i.e.,
\[
E(j_1) + q\cdot \beta(j_1) = E(j_2) + q\cdot \beta(j_2)
\]
\end{theorem}

\end{frame}


%------------------------------------------------

\begin{frame}\frametitle{Weak irrelevance result: Modigliani--Miller (1958)}

\begin{small}
\begin{proof}
Assume there exists an agent~$i$ having shares of both firms, i.e.,
\[
\alpha^i(j_1) >0 \et \alpha^i(j_2) >0
\]
Then there exists a family of state price deflators $(\lambda_s)_{s\in S}$ in $\R^S_{++}$ such that
\[
q = \sum_{s\in S} \lambda_s \kappa_s
\]
and
\[
\forall j\in \{j_1,j_2\},\quad E(j) = \sum_{s\in S} \lambda_s \delta_s(j)
\]
This implies that
\[
\forall j\in \{j_1,j_2\},\quad V(j) \equiv E(j) + q \cdot \beta(j) = \sum_{s\in S} \lambda_s y_s
\]
\end{proof}
\end{small}


\end{frame}


%------------------------------------------------

\begin{frame}\frametitle{Weak irrelevance result: Modigliani--Miller (1958)}

\textcolor{red}{What about if no agent has shares of both firms?}

\begin{proof}
\[
\ldots
\]
\end{proof}

\end{frame}

%------------------------------------------------

\begin{frame}\frametitle{Strong irrelevance result: Stiglitz (1969, 1974)}
\begin{theorem}
Consider a stock market equilibrium
\[
\left\{ (y(j),\beta(j))_{j\in J}, (q,E), (c^i,\theta^i,\alpha^i)_{i\in I}\right\}
\]
For any other financial plan $(\widehat{\beta}(j))_{j\in J}$ there exists another stock market equilibrium such that firms' values remain unchanged
\end{theorem}

\end{frame}

%------------------------------------------------

\begin{frame}\frametitle{Strong irrelevance result: Stiglitz (1969, 1974)}
More precisely, the following family
\[
\left\{ (y(j),\widehat{\beta}(j))_{j\in J}, (q,\widehat{E}), (c^i,\widehat{\theta}^i,\alpha^i)_{i\in I}\right\}
\]
is a stock market equilibrium where
\[
\widehat{E}(j) = E(j) + q \cdot [\beta(j) - \widehat{\beta}(j)]
\]
\[
\widehat{\theta}^i = \theta^i + \sum_{j\in J} \alpha^i(j) [\widehat{\beta}(j) - \beta(j)]
\]
\end{frame}

%------------------------------------------------

\begin{frame}\frametitle{Conditions for the validity of irrelevance results}

\begin{itemize}
\item Frictionless capital markets:
    \begin{itemize}
    \item no transaction costs
    \item no institutional restrictions on security trades
    \end{itemize}
\item No taxes
\item Investors and firms have access to the same security (bond) markets
\item Unlimited liability of shareholders
\item Investor do not try to infer information about the likelihood of states of nature by observing firm's financial policy
\end{itemize}
\end{frame}

%------------------------------------------------

\begin{frame}\frametitle{Conditions for the validity of irrelevance results}

\begin{itemize}
\item Linear effect: the change in the return on equity, induced by a modification of the firm's capital structure, is a linear combination of the returns on
the assets traded by the firm
\item Consider there are derivative securities (like options) whose payoff is a function (possibly non-linear) of the future value of another asset
\item When the firm's capital structure changes, not only the dividend of the stock changes but also that of the derivative securities written on the firm's
shares
\item In general, the firm's valuation will not be independent of the firm's capital structure
\end{itemize}

\begin{thebibliography}{8}


  \beamertemplatearticlebibitems

    \bibitem{Gottardi95}
    Gottardi, P.
    \newblock An analysis of the conditions for the validity of Modigliani--Miller theorem with incomplete markets.
    \newblock {\em Economic Theory} 5, 1995



    \end{thebibliography}

\end{frame}


%------------------------------------------------

\begin{frame}\frametitle{The role of taxes}

\begin{itemize}
\item Taxes are paid at $t=1$
\item We do not explain (model) the endogenous formation of tax rates
\item We do not explain the purpose of taxes (intermediaries, public projects)
\end{itemize}

\begin{block}{Three tax rates}
\begin{itemize}
\item corporate tax rate $t_c$ on corporate income $\delta_1= y_1 - \kappa_1 \cdot \beta$
\item personal tax rate $t_s$ on equity income
\item personal tax rate $t_b$ on financial assets (bonds) income
\end{itemize}
\end{block}

\begin{block}{Assumption}
\begin{itemize}
\item all tax payers face the same tax rates
\item tax rates are constant and independent of the level of income to which they apply
\end{itemize}
\end{block}

\end{frame}

%------------------------------------------------

\begin{frame}\frametitle{Assumptions}

Since taxes apply to income, we assume that
\begin{itemize}
\item every security $a$ is claim to a non-negative amount, i.e.,
\[
\forall a \in A, \quad \forall s\in S, \quad \kappa_s(a) \geq 0
\]
\item every firm $j$ never takes the risk to be bankrupt, i.e., the production-finance plan $(y,\beta)$ belong to $D$ where
\[
D \equiv \left\{ (y,\beta) \in Y \times \R^A \ \colon \ \forall s\in S, \quad y_s \geq \kappa_s \cdot \beta \right\}
\]
%\item investor can only buy securities
\end{itemize}

\end{frame}

%------------------------------------------------

\begin{frame}\frametitle{Investors' budget sets}

Each investor $i$ takes as given:
    \begin{itemize}
    \item each firm $j$'s production-finance plan $(y(j),\beta(j))$
    \item the security price vector $q$ and the equity price vector $E$
    \end{itemize}
and chooses
    \begin{itemize}
    \item a consumption plan $(c_0,c_1) \in \R_+ \times \R^S_+$
    \item a security portfolio $\theta\in \R^A_+$ and shareholdings $\alpha \in \R^J_+$
    \end{itemize}
satisfying budget restrictions

\end{frame}

%------------------------------------------------

\begin{frame}\frametitle{Investors' budget sets}

\begin{itemize}
\item solvency restriction at $t=0$
\begin{equation}\label{eq:budget-0-bis}
c_0 + q\cdot \theta + E \cdot \alpha \leq \omega^i_0 + (E + \delta_0) \cdot \xi^i
\end{equation}
\item solvency restriction at every $s\in S$ at $t=1$
\begin{equation}\label{eq:budget-s-bis}
c_s \leq \omega^i_s + (1-\tau_b)\kappa_s \cdot \theta + (1-\tau_s)(1-\tau_c)\delta_s \cdot \alpha
\end{equation}
\end{itemize}

The set of actions $(c,\theta,\alpha)$ satisfying (\ref{eq:budget-0-bis}) and (\ref{eq:budget-s-bis}) is denoted by $B^i(q,E,\tau)$

\end{frame}

%------------------------------------------------

\begin{frame}\frametitle{Stock market Equilibrium}

A stock market equilibrium with taxes is a list
\[
\left\{ (y(j),\beta(j))_{j\in J}, (q,E), (c^i,\theta^i,\alpha^i)_{i\in I}\right\}
\]
of
\begin{itemize}
\item production-finance plan $(y(j),\beta(j))$ for each firm $j$
\item security and equity prices $(q,E)$
\item action $(c^i,\theta^i,\alpha^i)$ of each investor $i$
\end{itemize}
such that
\end{frame}

%------------------------------------------------

\begin{frame}\frametitle{Stock market Equilibrium}

\begin{itemize}
\item[(a)] investor $i$'s action is optimal, i.e.,
\[
(c^i,\theta^i,\alpha^i) \in \argmax \{ U^i(c) \ \colon \ c \in B^i(q, E, \tau) \}
\]
\item[(b)] security and stock markets clear, i.e.,
\[
\forall a\in A, \quad \sum_{j\in J} \beta(j,a) = \sum_{i\in I} \theta^i(a) \et \forall j\in J, \quad \sum_{i\in I} \alpha^i(j) =1
\]
\item[(c)] commodity markets clear, i.e.,
\begin{eqnarray*}
c^p + \sum_{i\in I} c^i & = & \sum_{i\in I} \omega^i + \sum_{j\in J} y_s(j)
\end{eqnarray*}
where $c^p_0=0$ and
\[
c^p_s \equiv \sum_{j\in J} \left\{\tau_c + (1-\tau_c)\tau_s\right\} y_s(j) + \{\tau_b - \tau_c -(1-\tau_c)\tau_s\} \kappa_s \cdot \beta(j)
\]
\end{itemize}

\end{frame}

%------------------------------------------------

\begin{frame}\frametitle{Gains from corporate leverage}

Assume that there exist two firms $j_u$ and $j_\ell$ such that
\begin{itemize}
\item both firms have the same production plan, i.e., there exists $y\in Y(j_u) \cap Y(j_\ell)$ such that
\[
y= y(j_u) = y(j_\ell)
\]
\item firm $j_u$ is unlevered, i.e., $\beta(j_u) = 0$
\item firm $j_\ell$ is levered, i.e., $\beta(j_\ell)>0$
\item firm $j_\ell$ is never bankrupt, i.e.,
\[
\forall s\in S, \quad y_s \geq \kappa_s \cdot \beta(j_\ell)
\]
\end{itemize}


\end{frame}

%------------------------------------------------

\begin{frame}\frametitle{Gains from corporate leverage}

\begin{proposition}
Assume that there is a stock market equilibrium
\[
\left\{ (y(j),\beta(j))_{j\in J}, (q,E), (c^i,\theta^i,\alpha^i)_{i\in I}\right\}
\]
%with $\theta^i(a)>0$ for every security $a$ and every agent $i$.
Then
\[
V(j_\ell) = V(j_u) + \left[ 1 - \frac{(1-\tau_c)(1-\tau_s)}{1-\tau_b}\right] D(j_\ell)
\]
\end{proposition}

\end{frame}

%------------------------------------------------

\begin{frame}\frametitle{Role of corporate taxes}

Assume that personal income taxes coincide, i.e.,
\[
\tau_s = \tau_b
\]
then
\[
V(j_\ell) = V(j_u) + \tau_c D(j)
\]
\begin{itemize}
\item Firms have an incentive to raise capital via bond issues until the bankruptcy limit
\item Does not fit with data
%\item \alert{With taxes, do we still have that the correct objective (under complete markets) for firms is to maximize value?}
\end{itemize}

\end{frame}

%------------------------------------------------

\begin{frame}\frametitle{Role of personal tax system}

Assume that taxes matter, i.e.,
\[
(1-\tau_c)(1-\tau_s) \neq 1-\tau_b
\]
then
\begin{itemize}
\item if
\[
(1-\tau_c)(1-\tau_s) < 1-\tau_b
\]
then firms still have an incentive to raise capital via bond issues until the bankruptcy limit
\item if
\[
(1-\tau_c)(1-\tau_s) > 1-\tau_b
\]
then firms have an incentive to fully finance their capital through equity
\item these predictions do not fit with data
\end{itemize}

\end{frame}

%------------------------------------------------

\begin{frame}\frametitle{Do taxes may provide an optimal debt-to-equity ratio?}



\begin{block}{Theory}
\begin{thebibliography}{8}

  \beamertemplatearticlebibitems

    \bibitem{Miller77}
    Miller, M. H.
    \newblock Debt and taxes.
    \newblock {\em Journal of Finance} 32, 1977

    \end{thebibliography}
\end{block}

\begin{block}{Empirical Literature}


\begin{thebibliography}{8}

    \beamertemplatearticlebibitems

    \bibitem{Mackie-Mason90}
    Mackie-Mason, Jeffrey K.
    \newblock Do taxes affect corporate financing decisions?
    \newblock {\em Journal of Finance} 45, 1990

    \bibitem{DesayFoleyHines04}
    Desai, M. A., Foley, C. F., Hines Jr., J. R.
    \newblock A multinational perspective on capital structure choice and internal capital markets.
    \newblock {\em Journal of Finance} 59, 2004

    \end{thebibliography}
\end{block}

\end{frame}

%------------------------------------------------

\begin{frame}\frametitle{Limited liability and bankruptcy}

\begin{itemize}
\item Assume that only the riskless bond is available, i.e., $A=\{a_0\}$
\item Fix a firm and its production-finance plan $(y,Z)$
\item It may be the case that for some state $s$, we have
\[
\delta_s <0 \quad \textrm{or equivalently} \quad y_s < Z
\]
\item In order to protect (and attract) investors, shareholders have \alert{limited liability} in the sense that a new shareholder at $t=1$ is not called upon to
make payments when $\delta_s <0$
\end{itemize}
\end{frame}

%------------------------------------------------

\begin{frame}\frametitle{Limited liability and bankruptcy}

\begin{itemize}
\item If $\delta_s <0$ the firm is said to be {\em bankrupt}:
    \begin{itemize}
    \item all the firm's production $y_s$ is used to pay bondholders
    \item shareholders receive nothing
    \end{itemize}
\item It is as if the firm is issuing a bond which yields
\[
\forall s \in S, \quad b_s = \min\{y_s,Z\}
\]
\item One share of the firm yields
\[
\forall s\in S, \quad \widehat{\delta}_s = \max\{\delta_s,0\} = \max\{y_s - Z,0\}
\]
\end{itemize}

\end{frame}


%------------------------------------------------

\begin{frame}\frametitle{What are the markets available to investors}

\begin{block}{Irrelevance of the financial policy}
Every investor can purchase or sell any bond issued by firms
\end{block}

\begin{block}{Possible relevance of the financial policy}
\begin{enumerate}
\item Investors can purchase or sell only a ``pooled bond'' based on firms' bonds
\item Investors can only purchase bonds issued by firms
\end{enumerate}
\end{block}

\end{frame}
%------------------------------------------------

\begin{frame}\frametitle{Empirical Literature}


\begin{thebibliography}{8}

    \beamertemplatearticlebibitems

    \bibitem{SmithWarner79}
    Smith, C. and Warner, J.
    \newblock On financial contracting: an analysis of bond covenants.
    \newblock {\em Journal of Financial Economics} 7, 1979

    \bibitem{Weiss90}
    Weiss, L. A.
    \newblock Bankruptcy resolution: Direct costs and violation of priority of claims.
    \newblock {\em Journal of Financial Economics} 27, 1990


    \bibitem{OplerTitman94}
    Opler, T. C. and Titman, S.
    \newblock Financial distress and corporate performance.
    \newblock {\em Journal of Finance} 49, 1994



    \end{thebibliography}


\end{frame}

%------------------------------------------------

\begin{frame}\frametitle{Firms objectives}

\begin{itemize}
\item We focused our attention on the effects of the firm's decisions on the assets' rates of return
\item Other factors may guide the firms' financial decisions. The capital structure may be viewed as
    \begin{itemize}
    \item a signaling device: Ross~(1977), Leland and Pyle~(1977)
    \item a way of reducing agency costs: Jensen and Meckling~(1976)
    \item a way to allocate control rights between various claimants: Aghion and Bolton~(1988), Harris and Raviv~(1989)
    \end{itemize}
\end{itemize}

\end{frame}


%------------------------------------------------

\begin{frame}\frametitle{Irrelevance of the payout policy: A simple multi-period model}

\begin{itemize}
\item A multi period stock market: $t \in \{0,1,\ldots,T\}$
\item There is no uncertainty at every period
\item Bonds (or financial assets) are \alert{long-lived} and pay dividends at every period $t$
\item There is only one firm
\end{itemize}

\end{frame}

%------------------------------------------------

\begin{frame}\frametitle{Payout policies}
The firm chooses:
    \begin{itemize}
    \item An investment-production plan $(x,y) \in Z$, where
    \[
    x=(x_0,x_1,\ldots,x_{T-1}) \in \R_+^T
    \]
    and
    \[
    y=(y_1,\ldots,y_{T}) \in \R_+^T
    \]
    \item $x_t$ represents the units of the good invested (inputs) at period $t$
    \item $y_{t+1}$ represents the units of the good produced (outputs) at period $t+1$
    \item The technological restrictions are represented by the set $Z$
    \end{itemize}

\end{frame}

%------------------------------------------------

\begin{frame}\frametitle{Payout policies}
The firm chooses:
    \begin{itemize}
    \item A financial plan $\beta=(\beta_0,\beta_1,\ldots,\beta_{T-1})$ where
    \[
    \beta_t \in \R^A
    \]
    \item A dividend policy $n=(n_0,n_1,\ldots,n_{T-1})$ where $n_t$ represents the number of shares at
        \begin{itemize}
        \item the end of period~$t$, after payment of dividends $n_{t-1} \delta_t$
        \item beginning of period~$t+1$ before payments of dividends $n_t \delta_{t+1}$
        \end{itemize}
    \end{itemize}


\end{frame}

%------------------------------------------------

\begin{frame}\frametitle{Interpretation and notations}

Denote by $p_t$ the \alert{ex-dividend} price of one stock of firm $j$ at period $t$ \alert{after announcing the payout policy}

\begin{itemize}
\item If $n_t > n_{t-1}$ then firm $j$ is issuing shares and obtains the revenue $$p_t [n_t - n_{t-1}]$$
\item If $n_t < n_{t-1}$ then firm $j$ repurchases shares and spends $$-p_t [n_t - n_{t-1}]$$
\end{itemize}
\end{frame}

%------------------------------------------------

\begin{frame}\frametitle{Interpretation and notations}

Consider the firm's production and payout policy
\[
(x,y,\beta,n)
\]
\begin{itemize}
\item each investor $i$ chooses the stockholding $\alpha_t \geq 0$ \alert{in units of stocks} of the firm
\item one stock of the firm purchased at date $t$ is a claim to the dividend $\delta_{t+1}$ paid at period $t+1$
\[
n_t \delta_{t+1}=y_{t+1}- x_{t+1} - \kappa_{t+1}\beta_t  + p_{t+1} [n_{t+1} - n_{t}] + q_{t+1} [\beta_{t+1} - \beta_t]
\]
\item for notational convenience, we let
\[
y_0 =0, \quad x_T=0, \quad n_{-1}=1 \et \beta_{-1}=0
\]
\end{itemize}
\end{frame}

%------------------------------------------------

\begin{frame}\frametitle{Investors' budget sets}

Each investor $i$ takes as given:
    \begin{itemize}
    \item the firm's production-investment plan $(y,x)$ and payout policy $(\beta,n)$
    \item the securities' prices $q=(q_t)_t$ and the stock prices $p=(p_t)_t$
    \end{itemize}
and chooses
    \begin{itemize}
    \item a consumption plan $c=(c_t)_t \in \R^{T+1}_+$
    \item a security portfolio $\theta\in \R^{A\times T}$ and a stock-holdings vector $\alpha \in \R^{T}_+$
    \end{itemize}
satisfying ...

\end{frame}

%------------------------------------------------

\begin{frame}\frametitle{Investors' budget sets}
initially at $t=0$
\begin{equation}\label{eq:budget-0-ter}
c_0 + q_0\cdot \theta_0 + p_0 \cdot \alpha_0 \leq \omega^i_0 + (p_0 + \delta_0) \cdot \xi^i
\end{equation}
for each $t\in \{1,\ldots,T-1\}$,
\begin{equation}\label{eq:budget-t-ter}
c_t + q_t\cdot \theta_t + p_t \cdot \alpha_t \leq \omega^i_t + (\kappa_t + q_t) \cdot \theta_{t-1} + (\delta_t + p_t) \cdot \alpha_t
\end{equation}
and finally
\[
c_T \leq \omega^i_T + \kappa_T \cdot \theta_{T-1} + \delta_T \cdot \alpha_{T-1}
\]
The set of actions $(c,\theta,\alpha)$ satisfying the budget constraints is denoted by $B^i(q,p)$
\end{frame}

%------------------------------------------------

\begin{frame}\frametitle{Investors' budget sets}

initially at $t=0$
\begin{equation}\label{eq:budget-0-quatro}
c_0 + q_0\cdot (\theta_0 - \xi^i) + p_0 \cdot (\alpha_0 - 0) \leq \omega^i_0 +  \delta_0 \cdot \xi^i
\end{equation}
for each $t\in \{1,\ldots,T-1\}$,
\begin{equation}\label{eq:budget-t-quatro}
c_t + q_t\cdot (\theta_t - \theta_{t-1}) + p_t \cdot (\alpha_t-\alpha_{t-1}) \leq \omega^i_t + \kappa_t \cdot \theta_{t-1} + \delta_t  \cdot \alpha_{t-1}
\end{equation}
and finally
\[
c_T \leq \omega^i_T + \kappa_T \cdot \theta_{T-1} + \delta_T \cdot \alpha_{T-1}
\]
The set of actions $(c,\theta,\alpha)$ satisfying the budget constraints is denoted by $B^i(q,p)$
\end{frame}

%------------------------------------------------

\begin{frame}\frametitle{Stock market Equilibrium}

Given the firm's capital budgeting and payout policy
\[
(x,y,\beta,n)
\]
a stock market equilibrium is a list of
\[
\left\{ (q,p), (c^i,\theta^i,\alpha^i)_{i\in I}\right\}
\]
of
\begin{itemize}
\item security and share prices $(q,p)$
\item action $(c^i,\theta^i,\alpha^i)$ of each investor $i$
\end{itemize}
such that
\end{frame}

%------------------------------------------------

\begin{frame}\frametitle{Stock market Equilibrium}

\begin{itemize}
\item[(a)] investor $i$'s action is optimal, i.e.,
\[
(c^i,\theta^i,\alpha^i) \in \argmax \{ U^i(c) \ \colon \ (c,\theta,\alpha) \in B^i(q, p) \}
\]
\item[(b)] security and stock markets clear, i.e., for every $0\leq t<T-1$
\[
\forall a\in A, \quad \beta_t(a) = \sum_{i\in I} \theta^i_t(a) \et \sum_{i\in I} \alpha^i_t =n_t
\]
\item[(c)] commodity markets clear, i.e.,
\[
\sum_{i\in I} c^i = [y - x] + \sum_{i\in I} \omega^i
\]
\end{itemize}

\end{frame}

%------------------------------------------------

\begin{frame}\frametitle{No-arbitrage and state price deflator}

The definition is the same apart from stock markets clearing condition:
\[
\forall t \in \{0,\ldots,T-1\}, \quad \sum_{i\in I} \alpha^i_t = n_t
\]
The value $V_t$ of the firm at period $t$ is defined by
\begin{eqnarray*}
V_t & = & \underbrace{p_t n_t}_{E_t}  +  \underbrace{q_t\cdot \beta_t}_{D_t} + (y_t-x_t) %\underbrace{y_t - x_t}_{\DIV_t}
\end{eqnarray*}
\end{frame}

%------------------------------------------------

\begin{frame}\frametitle{Irrelevance of the payout policy}

By no-arbitrage, there exists a vector of state prices
\[
\lambda = (\lambda_0,\ldots,\lambda_T) \in \R^{T+1}_{++}
\]
such that for every $t \in \{0,\ldots,T-1\}$
\[
\lambda_t q_t = \lambda_{t+1} (\kappa_{t+1} + q_{t+1}) \et \lambda_t p_t = \lambda_{t+1} (p_{t+1} + \delta_{t+1})
\]
Observe that
\[
\lambda_0 V_0 = \sum_{t=0}^T \lambda_t (y_t-x_t)
\]
and
\[
\lambda_t V_t = \sum_{\tau=t}^T \lambda_\tau (y_\tau-x_\tau)
\]
The values of the firm are independent of the payout policy

\end{frame}

%------------------------------------------------

\begin{frame}\frametitle{How to explain the existence of an optimal capital structure?}

\begin{block}{Costs of agency}
These costs are due to conflict of interests among agents involved in the financial decisions:
\begin{itemize}
\item debtholders, new shareholders, old shareholders and managers all enter into negotiations for different reasons
\item bringing them into agreement is costly, concessions are required to achieve at least a second-best solution
\end{itemize}
\end{block}

\end{frame}

%------------------------------------------------

\begin{frame}\frametitle{How to explain the existence of an optimal capital structure?}

\begin{block}{Costs of asymmetric information}
Managers (insiders) may have a private information about the investment opportunities of the firm and related cash flows
\begin{itemize}
\item the information asymmetry may lead to some inefficiencies in the financing decisions of the firms
\item managers may use their choice of capital structure as a credible signal to the market of their private information
\end{itemize}
\end{block}

\end{frame}

\end{document}